\section{Life-facing synthetic data and a biological protocol}
\label{supp:life}

The carrier-side form of E1--E16 makes living organization a primary domain
of application, but substrate independence does not by itself turn a
biological analogy into a theorem instance.  This section therefore keeps
three evidential levels separate.  First, two computational
substrates supply finite evidence, with matched controls, for life-like maintenance,
drive separation, route dependence, and controlled higher-level readout
\citep{Tsiokos2026Life}.  Second, those measurements are compared with the
recognition conditions of the E-laws, with every missing bridge stated
explicitly.  Third, a concrete planarian regeneration study is specified as
an auditable biological protocol.  That protocol has not yet been run and is
not reported as evidence for any theorem.
The numbers below are those reported in the life-focused development
\citep{Tsiokos2026Life}; no result of the main paper depends on them.

\begin{table}[htbp]
\centering
\small
\begin{tabularx}{\textwidth}{@{}lXX@{}}
\toprule
Material & Evidence status & Admissible use here \\
\midrule
Particle and neural substrates &
Synthetic experiments with specified criteria and matched controls &
Finite evidence for the reported operational signals; not evidence that the
substrates are alive and not, without a bridge record, instances of an E-law
\\
E-law correspondence below &
Interpretive recognition map &
Identifies which measured fields could populate an E-law bridge and which
typed hypotheses remain unverified \\
HoloLife-1 planarian protocol &
Proposed and unrun &
Preregistered biological target, positive criteria, and falsifiers; no
positive biological result is claimed \\
\bottomrule
\end{tabularx}
\caption{Evidence-status discipline for the life-facing material.}
\label{tab:life-evidence-status}
\end{table}

% -----------------------------------------------------------------------------
\subsection{Two finite synthetic substrates}

The life-focused study instantiates two finite computational
substrates rather than treating ``life'' as a decorative interpretation
\citep{Tsiokos2026Life}.  The particle substrate has positions on a
two-dimensional torus, writable bonds $w_{ij}$, bounded counters $a_i$ and
$n_i$, a finite packaging field $S$, and optional operator-token and
meta-layer fields.  Its principal lenses expose the packaging field,
operator/interface mismatch, and windowed currents, affinities, motifs, and
protocol observables.  The neural/meta-layer substrate is a layered finite
lattice with bounded site variables, auxiliary fields, within-layer tokens
$W$ under a global budget, and inter-layer tokens $K$ under a per-site
budget.  Its lenses expose stroboscopic currents, cross-layer mismatch,
motif inventories, dictionaries, and intervention-conditioned decoding.

In both substrates, packaging is organized by the already specified
dynamics-induced completion template
\[
E_{\tau,f}(\nu)=Q_f\!\left((U_f\nu)P^\tau\right).
\]
The experiments do not compute the associated idempotence defect.  Their
stability claims are therefore operational maintenance and retention claims,
not measurements of exact closure.  Some neural ``snapshot'' observables
also depend on an entropy-production ledger; these are properly lenses on an
extended diagnostic state rather than unqualified observables of the
physical lattice state.

For the particle substrate, accepted $+1$ and $-1$ updates of each writable
coordinate type $y\in\{w,a,n,S\}$ give
\[
J_y=\frac{N_y^+-N_y^-}{T_w},
\qquad
A_y=\log\frac{N_y^++1}{N_y^-+1},
\qquad
\Sigma_{\mathrm{mem}}=\sum_y J_yA_y.
\]
The reported entropy-production quantities in both substrates are likewise
accepted-move proxies.  They omit rejected moves and, under the documented
conventions, proposal-ratio terms.  They are not exact path-space
Kullback--Leibler audits, need not be nonnegative in finite windows, and are
used only through null calibration and matched-regime separation.

\begin{table}[htbp]
\centering
\small
\begin{tabularx}{\textwidth}{@{}lXX@{}}
\toprule
Primitive & Particle realization & Neural/meta-layer realization \\
\midrule
P1 & writable bonds $w_{ij}$ & local $W$- and $K$-token exchange kernels \\
P2 & bounded apparatus counters $a_i$ & conserved token budgets and economy
diffusion \\
P3 & explicit protocol phase/clock & deterministic kernel cycle keyed to
step index \\
P4 & quantized counters $n_i$ & template-flip kernel in the later/meta
presets \\
P5 & writable packaging field $S$ & barrier/closure field and structured
token-mediated packaging \\
P6 & drive toggle and $\mu$-contexts & nonconservative work term with a
separate drive toggle \\
\bottomrule
\end{tabularx}
\caption{Realizations of the primitive roles in the two synthetic substrates.  A
primitive being implemented does not establish an E-law whose setup uses
that primitive.}
\label{tab:life-primitive-map}
\end{table}

\subsubsection{Null calibration and drive separation}

The audit-first sequence begins with P3 and P6 disabled.  The relevant proxy
audits are near zero while the dynamics remain nondegenerate
(\cref{tab:life-null}).  Only after this calibration is P6 enabled under a
matched regime.  The resulting separation is reported in
\cref{tab:life-drive}.  These are measurements of a controlled drive channel;
they do not establish biochemical metabolism, thermodynamic autonomy, or an
exact arrow-of-time functional.

\begin{table}[htbp]
\centering
\small
\begin{tabularx}{\textwidth}{@{}llX@{}}
\toprule
Substrate & Null regime & Exact reported proxy summary \\
\midrule
Particles & P3 off, P6 off &
$\Sigma_{\mathrm{mem}}=\texttt{sigmaMem}\approx3\times10^{-4}$;
M6 motif affinities $=0$; accepted-move EP proxy
$\approx3.5\times10^{-6}$ with confidence-interval half-width
$\approx4.7\times10^{-5}$, containing zero \\
Neural & P3 off, P6 off &
Specified bound
$|\texttt{epMicroRateWindowLast}|\le2\times10^{-4}$ for the specified seeds;
at $24\times24$ over three seeds, mean EP
$\approx-3.9\times10^{-6}$, maximum confidence-interval half-width
$\approx2.8\times10^{-4}$, and acceptance $\approx0.54$ \\
\bottomrule
\end{tabularx}
\caption{Executed null-regime calibration.  The sign of the finite-window
neural proxy is not assigned thermodynamic meaning.}
\label{tab:life-null}
\end{table}

\begin{table}[htbp]
\centering
\small
\begin{tabularx}{\textwidth}{@{}llX@{}}
\toprule
Substrate & Regime & Exact reported separation \\
\midrule
Particles & Null, P6 off &
$\texttt{sigmaMem}\approx3\times10^{-4}$; M6 affinities $=0$; EP proxy
$\approx3.5\times10^{-6}$ with a confidence interval containing zero \\
Particles & Drive, P6 on &
$\texttt{sigmaMem}\approx2.1\times10^{-3}$;
$\texttt{M6W}\approx0.814$,
$\texttt{M6N}\approx0.374$,
$\texttt{M6A}\approx0.392$; EP proxy $\approx4.26\times10^{-3}$ \\
Neural & Null, P6 off &
$|\texttt{null\_mean}|\le5\times10^{-4}$ with nondegenerate acceptance \\
Neural & Drive, P6 on &
$\texttt{drive\_mean}\ge10^{-3}$ and
$\texttt{drive\_mean}-\texttt{null\_mean}\ge5\times10^{-4}$ at
$B_k=2$ and $\eta_{\mathrm{drive}}=2$ \\
\bottomrule
\end{tabularx}
\caption{Executed P6 null--drive separation with P3 held off.}
\label{tab:life-drive}
\end{table}

\subsubsection{Protocol effect without a directionality overread}

The neural P3 experiment compares a deterministic kernel schedule with a
matched control at one observation lens.  The schedule changes the
stroboscopic-current statistic in every reported row
(\cref{tab:life-protocol}), but the protocol is keyed to the step index.
Unless protocol phase is included in the state, the observed process is
time-inhomogeneous.  Accordingly, the measurement is evidence of a
route/protocol effect and is not used as a directionality certificate.

\begin{table}[htbp]
\centering
\scriptsize
\begin{tabular}{@{}lccccc@{}}
\toprule
Setting & Control L2 & Protocol L2 & Relative change & Accept (control) &
Accept (protocol) \\
\midrule
seed 1, $\eta=0.5$ (additional) & 0.0349 & 0.0262 & 0.250 & -- & -- \\
seed 1, $\eta=1.0$ (additional) & 0.0379 & 0.0271 & 0.284 & -- & -- \\
seed 1, $\eta=2.0$ (specified) & 0.0311 & 0.0255 & 0.179 & -- & -- \\
\addlinespace
seed 1, $\eta=1.0$ (replication) & 0.0379 & 0.0271 & 0.284 & 0.0103 & 0.0107 \\
seed 2, $\eta=1.0$ (replication) & 0.0326 & 0.0280 & 0.141 & 0.0116 & 0.0107 \\
seed 3, $\eta=1.0$ (replication) & 0.0313 & 0.0280 & 0.105 & 0.0107 & 0.0108 \\
\bottomrule
\end{tabular}
\caption{Executed neural protocol diagnostic.  L2 is the reported
\texttt{strobe\_current\_l2\_window}.  The $\eta=2$ row is the specified
case; the other rows are additional settings.}
\label{tab:life-protocol}
\end{table}

\subsubsection{Maintenance and viability probes}

The maintenance experiments ask whether a finite system preserves or
recovers a structured target under noise, traversal, hazard, or deadline.
The targets, hazards, and deadlines are imposed by the experimental design.  Thus
the results establish controlled maintenance competencies, not a free
metabolism and not an autonomous origin of the repair objective.

\begin{table}[htbp]
\centering
\scriptsize
\begin{tabularx}{\textwidth}{@{}llX@{}}
\toprule
Substrate & Task & Reported outcome \\
\midrule
Particles & Code maintenance &
Null: $\texttt{err}=0.43\pm0.09$,
$\texttt{Sdiff}=6.86\pm0.25$, $\texttt{epRate}=0$.
Drive: $\texttt{err}=0$,
$\texttt{Sdiff}=0.83\pm0.06$, $\texttt{epRate}=0.016$ \\
Particles & Gated traversal &
Ungated: $\texttt{err}=0.23$, $\texttt{Sdiff}=0.49$.
Gated moving: $\texttt{err}=0$, $\texttt{Sdiff}=0$.
Gated static: $\texttt{Sdiff}=2.58$ and no recovery \\
Particles & Deadline selection &
$\Delta\texttt{uptimeTail}=0.285$ with confidence interval
$[0.14,0.43]$; clock cost
$\Delta\texttt{EPClockRate}=0.083$ \\
Neural & Paired hazard response &
$\texttt{spike\_paired}=0.063$,
$\texttt{realloc\_paired}=0.015$,
$\texttt{recovery\_paired}=0.060$; declared reallocation threshold
$0.002$ \\
\bottomrule
\end{tabularx}
\caption{Executed maintenance and viability probes.  The neural hazard
response is measured with paired runs.}
\label{tab:life-maintenance}
\end{table}

\subsubsection{Refined-lens motifs and controlled readout}

The neural substrate also exposes motif histograms and motif-transition
statistics at a refined diagnostic lens.  The inventory experiment requires coverage at
least $0.60$, Jensen--Shannon divergence at least $0.01$, and propagation
score at least $0.02$; all three seeds exceed those thresholds.  The transition experiment reports
the transition divergence in the same finite regime.

\begin{table}[htbp]
\centering
\scriptsize
\begin{tabular}{@{}lcccccc@{}}
\toprule
Experiment & Seed & Coverage (pre) & Coverage (hazard) & JSD (pre/hazard) &
JSD (transition) & Propagation \\
\midrule
inventory & 1 & 1.0 & 1.0 & 0.0312 & -- & 0.644 \\
inventory & 2 & 1.0 & 1.0 & 0.0324 & -- & 0.587 \\
inventory & 3 & 1.0 & 1.0 & 0.0316 & -- & 0.609 \\
\addlinespace
transition & 1 & 1.0 & -- & 0.0312 & 0.0383 & 0.644 \\
transition & 2 & 1.0 & -- & 0.0324 & 0.0418 & 0.587 \\
transition & 3 & 1.0 & -- & 0.0316 & 0.0424 & 0.609 \\
\bottomrule
\end{tabular}
\caption{Executed motif-inventory and motif-transition results.}
\label{tab:life-motifs}
\end{table}

Four further experiments inject controlled labels and evaluate the resulting
motif-level statistics against a within-run shifted-label null.  Here
``semantics'' means intervention-conditioned discriminability and recovery
of injected labels, not linguistic meaning, awareness, or subjective
content.  The $p$-values in \cref{tab:life-intervention} are experiment-local
shift-null controls; no global family-wise error claim is made.

\begin{table}[htbp]
\centering
\scriptsize
\begin{tabular}{@{}lclccl@{}}
\toprule
Experiment & Seed & Metric & Value & $p$-value & Reported separation \\
\midrule
directional focus & 1 & focus delta & 1.478 & 0.005 & spike control 0.269 \\
directional focus & 2 & focus delta & 1.453 & 0.000 & spike control 0.225 \\
directional focus & 3 & focus delta & 1.521 & 0.000 & spike control 0.384 \\
\addlinespace
dictionary polarity & 1 & dictionary delta & 0.064 & 0.005 & JSD in/out 0.032 \\
dictionary polarity & 2 & dictionary delta & 0.108 & 0.005 & JSD in/out 0.060 \\
dictionary polarity & 3 & dictionary delta & 0.075 & 0.005 & JSD in/out 0.039 \\
\addlinespace
phrase alignment & 1 & alignment & 0.999 & 0.005 & JSD out/in 0.665 \\
phrase alignment & 2 & alignment & 1.000 & 0.005 & JSD out/in 0.671 \\
phrase alignment & 3 & alignment & 1.000 & 0.005 & JSD out/in 0.667 \\
\addlinespace
phrase decoding & 1 & token accuracy & 1.000 & 0.005 & JSD out/in 0.665 \\
phrase decoding & 2 & token accuracy & 1.000 & 0.005 & JSD out/in 0.671 \\
phrase decoding & 3 & token accuracy & 1.000 & 0.005 & JSD out/in 0.667 \\
\bottomrule
\end{tabular}
\caption{Executed intervention-conditioned readout and shifted-label-null
results.}
\label{tab:life-intervention}
\end{table}

Three other readout experiments failed their shifted-null criteria.  They count
against the readout claims, not for them.  In particular, the successful
intervention-conditioned decoding does not follow merely from the presence
of motifs.

% -----------------------------------------------------------------------------
\subsection{What the synthetic evidence does and does not instantiate}

The measurements above give life-facing substance to the abstract catalog,
but they are not silently promoted to E-law instances.  The recognition map
in \cref{tab:life-e-law-map} states the strongest source-faithful relation.

\begin{table}[htbp]
\centering
\scriptsize
\begin{tabularx}{\textwidth}{@{}lXX@{}}
\toprule
E-law surface & Existing life-facing signal & Missing bridge or control \\
\midrule
E1 Internalization and E3 Self-Maintaining Reclosure &
Driven code maintenance, gated traversal, perturbation recovery, and explicit
costs are finite repair/maintenance candidates &
The target and hazard are externally installed.  A full instance still
requires a carried defect record, carrier-produced source repair, the E1
closed-loop record, and the E3 apparatus reclosure and audit
chain \\
E5 Reclosure Collapse &
Null, ungated, and gated-static failures provide genuine separating controls
for maintenance claims &
Failure of an experimental task is not yet the typed exhaustion of moves,
budget, and reclosure machinery required by E5 \\
E6 Priced Exposure and E9 Strict Probe Acquisition &
The neural carrier has explicit $W/K$ budgets; the particle carrier exposes
finite diagnostic families &
No complete native/probe catalogue, strict-family test, $\Xi$ residual,
positive-cost ledger, or KKT allocation certificate links these runs to E6
or E9 \\
E7 Alarm and E8 Control-Price &
The paired neural hazard produces a spike, resource reallocation, and
recovery; the particle deadline run reports a benefit and an explicit proxy
cost &
The E7 disposition inventory and carried alarm record are not built, and the
E8 signal is not linked to a declared binding constraint, dual variable,
lineage audit, and capture controls \\
E10 Cognitive Demarcation &
Motifs, dictionaries, and intervention-conditioned decoding furnish
decodable correlates under controls &
Decodability is not cognition.  A positive E10 instance additionally needs
the same carrier's repair-and-gate evidence, carried decision link,
counterfactual ablations, and status partition \\
E14 Reconsolidation and E15 Offline Reclosure &
The controlled representations provide candidate
records to examine &
No retrieval-contingent record revision, canonical E14 status, online/offline
debt flow, exchange census, or duty-cycle theorem is measured here \\
E16 Adaptability &
The P3 study establishes a repeatable protocol-dependent change in one
stroboscopic statistic &
A protocol effect alone is not coherent adaptability.  E16 still requires
fixed support, matched current endpoints, exact endpoint-conditioned repair
distributions, E15 reconciliation, complete counterpart inventories, and
bounded-perturbation survival \\
\bottomrule
\end{tabularx}
\caption{Cautious recognition map from executed synthetic evidence to the
E-law surfaces.  Every row records both positive material and the missing
bridge.}
\label{tab:life-e-law-map}
\end{table}

No result above directly instantiates E4 compilation, E11 support rewrite,
E12 individuation, or E13 repair transport.  The two substrates do,
however, show why living organization belongs beside cognition and society
rather than underneath them: maintenance, resource allocation, route
dependence, and controlled readout are already nontrivial carrier-side
phenomena before any cognitive interpretation is attempted.

% -----------------------------------------------------------------------------
\subsection{HoloLife-1: a proposed biological instantiation}
\label{sec:hololife-protocol}

The synthetic evidence does not substitute for a biological test.  The
HoloLife-1 proposal stated here targets planarian regeneration, motivated by
work on bioelectric control and regenerative pattern memory
\citep{PezzuloLevin2015}.  Its central object is not a black-box predictor and
not a claim that the organism ``believes'' a target morphology.  It is a
finite, audit-led descent/repair claim: two histories that agree at a
declared current quotient may separate under a future recut challenge, and
the proposed study asks whether a carried bioelectric predicate and route
record lawfully repair that failure.

\paragraph{Finite carrier and audit ledger.}
Let $H=\{h_i\}$ be a frozen finite set of regeneration-history records.  Each
record contains
\[
\begin{aligned}
h=(&\text{source/study},\ \text{species},\ \text{fragment type},\\
&\text{amputation geometry},\ \text{treatment route},
\ \text{washout state},\\
&\text{time/stage},\ \text{current anatomy},
\ \text{current molecular markers},\ \text{bioelectric readout},\\
&\text{recut challenge},\ \text{future outcome},\ \text{reset
intervention},\ \text{audit ledger}).
\end{aligned}
\]
The audit ledger is carrier data.  It records instruments, declared
thresholds, visible and suppressed fields, empirical bridges, species and
protocol compatibility, remaining residuals, and any target coarsening.  In
particular, public molecular atlases may populate a declared bridge for the
current molecular quotient only when species, time, and protocol
compatibility are recorded; they may not be silently merged with the
perturbation histories.

\paragraph{Four quotients and one held-out challenge.}
The proposed quotients are
\[
Q_0=Q_A\vee Q_G,
\qquad
Q_B=Q_0\vee h_B,
\qquad
Q_R=Q_B\vee\rho,
\qquad
M=Q_\Theta.
\]
Here $Q_A$ records current visible anatomy, $Q_G$ the declared current
molecular/marker state, $h_B$ a finite bioelectric predicate, $\rho$ a
route-memory residue, and $Q_\Theta$ the future target-morphology quotient.
The held-out map $F_{\mathrm{recut}}:H_t\to H_{t+\Delta}$ is read into
\[
r_\Theta F_{\mathrm{recut}}(h)\in\{W,C,D,R,A\},
\]
where $W$ is a wild-type single-headed target, $C$ a cryptic altered target,
$D$ a double-headed target, $R$ a reset/repaired target, and $A$ an abnormal
or failed package.  These labels are proposed finite outcome classes; they
do not earn P5 objecthood merely by being assigned.

\paragraph{Test 1: split-pair obstruction.}
The primary target is the already specified descent obstruction
\[
\Omega(Q_0;F,r_\Theta)
=\{(h,h'):\ Q_0(h)=Q_0(h')\ \wedge\
r_\Theta F(h)\ne r_\Theta F(h')\}.
\]
The positive first-stage criterion is
$\Omega(Q_0;F,r_\Theta)\ne\varnothing$.  This says that the declared present
anatomy/marker quotient is inadequate for the future recut readout.  It is
stronger and more falsifiable than reporting predictive accuracy alone.

\paragraph{Test 2: strict bioelectric source repair.}
The candidate $h_B$ must first split at least one old fiber:
\[
h_B\notin\Sigma_{Q_0},
\qquad
I(h_B\mid Q_0)
=\frac1N\sum_{C\in H/Q_0}\min(n_{C,0},n_{C,1})>0.
\]
It must then reduce the preregistered obstruction,
\[
\widehat\omega(Q_B)\ll\widehat\omega(Q_0),
\qquad Q_B=Q_0\vee h_B,
\]
rather than merely improve a fitted score.  Entropy-matched random
predicates are explicit comparators.  Target coarsening is recorded as a
different repair: collapsing head/tail outcomes to ``viable'' cannot be
credited as an explanation of the original distinction.

\paragraph{Test 3: hiddenness under exposure.}
The HoloLife-1 proposal uses the empirical witness-pair frequency
below.  To avoid identifying it with D2's closure-deficit surplus without a
bridge, write it here as
\[
S_{\rm pair}(Q,M)=\Pr[h\sim_Q h'\ \wedge\ h\not\sim_M h'],
\]
with positive pattern
\[
S_{\rm pair}(Q_0,M)>0,
\qquad
S_{\rm pair}(Q_B,M)\ll S_{\rm pair}(Q_0,M).
\]
Threshold and treatment-strength sweeps test whether the reduction is tied
to exposure of the determining predicate rather than to an arbitrary large
effect.  A residual that remains after exposure stays statused rather than
being declared solved.  Reading it as an instance of D2 would require the
registered closure-deficit score or an explicit bridge from this finite pair
frequency to that score.

\paragraph{Test 4: route memory without the protocol trap.}
The comparison fixes two routes:
\[
\begin{aligned}
\gamma &: \text{untreated regeneration}
\to\text{current single-headed state}\to\text{recut},\\
\eta &: \text{transient bioelectric perturbation}
\to\text{washout}\to\text{current single-headed state}\to\text{recut}.
\end{aligned}
\]
The proposed holonomy-memory signature is
\[
Q_0(\gamma h)=Q_0(\eta h),
\qquad
M(\gamma h)\ne M(\eta h),
\]
with the route residue $\rho_{\gamma,\eta}$ added on the source side through
$Q_R=Q_B\vee\rho_{\gamma,\eta}$.  This would be a physiological formed
record, not a cognitive state.  It would not by itself establish P6 or
thermodynamic directionality.  Any P6 claim must separately lift treatment,
phase, and protocol into state and survive the corresponding audit.

\paragraph{Test 5: package formation.}
The proposed target states earn P5 status only through the existing finite
completion test
\[
E_{\tau,q}(\mu)=U_q\!\left(Q_q(\mu K^\tau)\right),
\qquad
\delta_{\tau,q}
=\frac12\max_z\sum_{z'}
\left|(E_{\tau,q}^2-E_{\tau,q})(z,z')\right|.
\]
The preregistered positive condition is
$\delta_{\tau,Q_\Theta}\ll1$ together with at least two nontrivial stable
packages.  A constant map, unstable classifier, or post hoc outcome bin does
not pass.

\paragraph{Test 6: lower-layer gating or rewrite.}
For target package $\theta$, the proposed lower-layer kernel is
\[
K_B^\theta(b,b')=
\Pr(B_{t+\Delta}=b'\mid B_t=b,\Theta_t=\theta,\phi_t),
\]
where $\phi_t$ includes phase, treatment, washout, recut protocol, and
instrument state.  The registered comparisons are kernel divergence,
support symmetric difference, cycle-rank change, and spectral-gap change.
A support or feasibility change is a P2 candidate; a change in reachable
classes, absorbing states, cycle rank, spectral behavior, or the induced
packaging map is a P1 candidate.  Small parameter drift alone is below the
registered threshold.  The distinctive positive target is a difference
among $K_B^W,K_B^C,K_B^D$ under the same lifted external protocol.

\paragraph{Test 7: residual and currency discharge.}
Let $L=\{Q_A,Q_G\}$ be the native current probes and let
$D=\{h_B,\rho_{\gamma,\eta},Q_\Theta\}$ be the proposed dissolving family.
The existing adequacy residual is
\[
\Xi_C(D\mid L)
=K_{DD}-K_{DL}K_{LL}^{\dagger}K_{LD}.
\]
The protocol requires the residual after adding the bioelectric and route
records to fall below a declared budget:
\[
\Xi_C(D_{\mathrm{remaining}}\mid L\vee h_B\vee\rho)
<\text{declared budget}.
\]
If it does not, the biological needle remains undisolved.  A later shadow
price may be interpreted only when the constrained budget problem, binding
lineage, and dual certificate are all actually supplied; a fitted
coefficient is not an E8 control price.

\paragraph{Role-status ledger.}
The proposal ends in the finite status table of
\cref{tab:hololife-role-ledger}, not in an untyped claim that all six
primitives are present.

\begin{table}[htbp]
\centering
\scriptsize
\begin{tabularx}{\textwidth}{@{}lXX@{}}
\toprule
Role & Positive witness sought & Demotion or failure \\
\midrule
P1 rewrite &
$K_B^\theta$ changes reachable classes, absorbing states, cycle rank,
spectral gap, or package map by target state &
Differences vanish after phase, protocol, treatment, and instrument lifting
\\
P2 gating &
Target package changes the support or feasibility of bioelectric/regenerative
transitions &
Only weak parameter drift; no support or feasibility change \\
P3 holonomy &
Routes agree at $Q_0$ but differ at $M$ &
Difference is explained by an unrecorded treatment, time, or protocol field
\\
P4 staging/forking &
$W,C,D,R,A$ form a typed branch tree with transition records and stage
indices &
Labels are post hoc outcome bins \\
P5 packaging &
$Q_\Theta$ has low idempotence defect and multiple stable packages &
Constant map or unstable classifier \\
P6 audit &
Thresholds, instruments, bridges, residuals, repairs, coarsenings, and
nonclaims are carried and statused &
Any silent bridge, silent translation, or unstatused residual \\
\bottomrule
\end{tabularx}
\caption{Proposed HoloLife-1 role-status ledger.  Every row is a test to be
run, not an observed biological result.}
\label{tab:hololife-role-ledger}
\end{table}

\subsubsection{Preregistered positive package}

A strong positive result requires the conjunction in
\cref{tab:hololife-positive}.  Passing only one or two rows would support a
narrower biological statement, not the proposed full-stack conclusion.

\begin{table}[htbp]
\centering
\small
\begin{tabularx}{\textwidth}{@{}lX@{}}
\toprule
Preregistered criterion & Licensed conclusion \\
\midrule
$\Omega(Q_0;F,r_\Theta)\ne\varnothing$ &
The declared present anatomy/marker quotient fails descent to future form \\
$h_B\notin\Sigma_{Q_0}$ and
$\widehat\omega(Q_B)\ll\widehat\omega(Q_0)$ &
The bioelectric predicate is a strict source extension that repairs, rather
than merely redescribes, the obstruction \\
$S_{\rm pair}(Q_B,M)\ll S_{\rm pair}(Q_0,M)$ &
Predictive hiddenness is reduced under the declared exposure \\
$Q_0(\gamma h)=Q_0(\eta h)$ but
$M(\gamma h)\ne M(\eta h)$ &
The declared routes require a source-side memory residue \\
$\delta_{\tau,Q_\Theta}\ll1$ with at least two nontrivial stable packages &
Target morphology passes the proposed P5 package test \\
Target-conditioned $K_B^\theta$ differ under the same lifted protocol &
The formed package is associated with lower-layer gating or rewrite, as
classified by the registered support and kernel tests \\
$\Xi_C(D_{\mathrm{remaining}}\mid L\vee h_B\vee\rho)$ is below the declared
budget &
The residual is discharged in the registered currency rather than hidden by
prose \\
All role channels are statused &
No primitive is credited through an unrecorded bridge or residual \\
\bottomrule
\end{tabularx}
\caption{Conjunctive positive criteria for the proposed biological
protocol.}
\label{tab:hololife-positive}
\end{table}

\subsubsection{Preregistered falsifiers}

The protocol is designed to make failure informative and easy to report.

\begin{table}[htbp]
\centering
\small
\begin{tabularx}{\textwidth}{@{}lX@{}}
\toprule
Observed failure & Consequence \\
\midrule
$Q_0$ already descends to future target morphology &
There is no split-pair obstruction and no hidden target package relative to
the declared current quotient \\
$h_B$ is definable from $Q_0$ &
There is no strict bioelectric source extension \\
An entropy-matched random predicate repairs descent as well as $h_B$ &
The proposed source repair lacks the registered biological specificity \\
Only target coarsening repairs the obstruction &
The distinction has been erased rather than explained \\
Route memory disappears after treatment, time, phase, and protocol are
lifted &
The apparent holonomy was a protocol or omitted-variable artifact \\
$\delta_{\tau,Q_\Theta}$ remains high or only a constant package survives &
Target morphology has not earned P5 objecthood \\
The adequacy residual remains above budget after $h_B$ and $\rho$ are added &
The proposed needle remains unresolved \\
$K_B^\theta$ does not vary with target package under the fixed lifted
protocol &
This evidence pack does not support the lower-layer rewrite hypothesis \\
Any bridge or residual cannot be statused &
The full Six-Birds biological package is rejected even if individual
correlations remain \\
\bottomrule
\end{tabularx}
\caption{Falsifiers for the proposed HoloLife-1 study.}
\label{tab:hololife-falsifiers}
\end{table}

\paragraph{Implementation boundary.}
The proposal calls for a frozen evidence pack containing perturbation route,
current morphology, future recut outcome, reset behavior, and declared
bioelectric readouts; digitized voltage-map features; and molecular atlases
used only through explicit compatibility bridges.  The core analysis is a
finite obstruction/repair/audit certificate using finite-state estimation,
uncertainty intervals, bootstrap or permutation controls, and
entropy-matched random predicates.  A simulator may be attached only as a
model-realization map.  It cannot replace the primary biological carrier or
convert this unrun protocol into present evidence.

\paragraph{Scope of a future positive result.}
Even the full conjunction would not define biological life, prove that all
living systems realize E1--E16, establish consciousness, or derive a
thermodynamic arrow from protocol holonomy.  It would provide a finite
biological realization of a narrower normal form: current-quotient failure,
strict physiological source repair, exposure-sensitive predictive surplus,
route-carried memory, package formation, lower-layer transition change, and
audited residual discharge.  Until the protocol is executed, that sentence
states the target of the experiment and nothing more.

% =============================================================================
